% !TEX root = ../../main.tex

\section{模型版本发布与识别服务联动实现}

模型版本发布与识别服务联动是 HearBridge 系统模型维护闭环中的关键流程。
后台管理端完成模型训练或上传模型文件后，需要在服务端登记模型版本信息，包括版本名称、模型文件地址、标签映射文件地址、支持动作数量、训练样本数量、准确率和发布状态等内容。
服务端通过模型版本表保存这些信息，并向后台管理端提供模型版本查询、发布和状态维护接口。

模型发布的核心目标是保证系统运行时使用明确的当前模型版本。
管理员在后台选择某一模型版本并发布后，服务端需要更新数据库中的发布状态，使同一时刻只有一个模型处于当前发布状态。
随后，服务端调用 Python 识别服务的模型重载接口，使实时识别和视频识别使用最新模型文件。

模型版本发布与识别服务联动流程如图~\ref{fig:model-release-flow} 所示。

\WhuImageFigure[0.78\textwidth][0.52\textheight]{figures/chapter05/model-release-flow.pdf}{模型版本发布与识别服务联动流程图}{fig:model-release-flow}

模型版本发布联动还需要考虑服务状态和错误处理。
如果 Python 识别服务加载模型失败，服务端应向后台管理端返回明确错误，而不能只更新数据库状态。
若模型文件地址、标签映射文件或版本参数不完整，也应在发布前进行校验。
通过这种方式，模型版本管理不仅是后台数据维护功能，也承担识别服务运行状态切换的控制作用。

在移动端实时识别场景中，识别结果中可以返回当前模型版本名称，使用户和开发者能够判断识别结果来自哪个模型版本。
该信息对于后续排查识别效果、分析模型更新前后差异具有一定价值。
通过服务端模型版本管理和 Python 识别服务加载机制的配合，系统形成了从样本维护、模型训练、版本登记到模型发布的基本闭环。
